\documentclass[10pt]{article}

\usepackage[utf8]{inputenc}

\usepackage[T1]{fontenc}

\usepackage{amsmath}

\usepackage{amsfonts}

\usepackage{amssymb}

\usepackage[version=4]{mhchem}

\usepackage{stmaryrd}

\usepackage{graphicx}

\usepackage[export]{adjustbox}

\graphicspath{ {./images/} }



\begin{document}

странства были отделимы посредством окрестностей, т. е. содержались в дизъюнктных открытых множествах (совокупность двух или более множеств дизъюнктна или состоит из дизъюнктных множеств, если эти множества попарно не имеют общих элементов). Аксиома отделимости Хаусдорфа иначе наз. а к с и ом о й $T_{2}$, а удовлетворяющие ей Т. п. наз. х а у с д о рф о в ы м и или $T_{2}$-П р о с т р а н с т в а м и. Определив әти пространства, Ф. Хаусдорф (F. Hausdorff) в 1914 впервые открыл достаточно широкий и в то же время достаточно богатый свойствами класс Т. П. и тем удовлетворил уже вполне назревшую к тому времени потребность математики. Дальнейшее развитие общей топологии исходит именно из хаусдорфовых пространств. Ослаблением аксиомы $T_{2}$ является а кс и о м а $T_{1}$, требующая, чтобы каждая из любых двух точек Т. П. имела окрестность, не содержащую вторую точку. Это требование равносильно требованию замкнутости в данном пространстве всякого множества, состоящего из конечного числа точек. Пространства, удовлетворяющие этому требованию, наз. $T_{1}$-П $\mathrm{p}$ о с тр а н с т в а м и. Еще более широкий класс Т. П. образуют $T_{0}$-П р о с т р а н с т в а, т. е. иространства, в к-рых выполнена а к с и о м а $T_{0}$ (а к с и о м а К о л м ог о р о в а): из любых двух точек этого пространства по крайней мере одна имеет окрестность, не содержащую другую точку. $T_{0}$-пространства могут состоять из конечного числа и даже из двух точек, если одна точка образует замкнутое множество, но не открытое, а другая, наоборот, образует открытое, но не замкнутое множество (с в я з н о е д в о е т о ч и е).



А к с и о м а $T_{3}$ требует, чтобы произвольная точка пространства и произвольное, не содержащее эту точку замкнутое множество были отделимы окрестностями, т. е. содержались соответственно в двух дизъюнктных открытых множествах. Пространства, удовлетворяющие аксиоме $T_{3}$, наз. $T_{3}$-п р о с т р а н с т в а м и. $T_{3}$-пространство может не удовлетворять аксиоме $T_{1}$, таково, напр., связное двоеточие. $T_{0}$-пространства, удовлетворяющие аксиоме $T_{1}$, наз. р е г у л я р н ы м и; все онихаусдорфовы пространства. Т. П. наз. $T_{4}$-П р о с т р а нс т в о м, если в нем всякие два дизъюнктные замкнутые множества имеют дизъюнктные окрестности. $T_{4}$ пространства, являющиеся в то же время $T_{1}$-пространствами, наз. н о р м а л ь н м и, все они регулярны и подавно хаусдорфовы. Всякое множество $X_{0}$ точек T. п. $X$, (\$) получает топологию, однозначно определенную топологией (j, и таким образом становится топологич. подпространством пространства $X$, (\$). При этом всякое подпространство пространства, удовлетворяющего аксиоме $T_{i}, \quad i=0,1,2,3$, также удовлетворяет аксиоме $T_{i}$, т. е. аксиомы $T_{i}$ наследуются всеми подпространствами данного пространства. Аксиома $T_{4}$ этим свойством не обладает: существуют нормальные пространства $X$, не все (даже открытые в $X$ ) подпространства к-рых нормальны. Однако если $X_{0}$ замкнутое множество в нормальном пространстве $X$, то подпространство $X_{0}$ нормально.



До сих пор отделимость точек и множеств понималась в смысле наличия у них дизъюнктных окрестностей. Однако в современной топологии имеет значение и т. н. функциональная отделимость, впервые введенная П. С. Урысоном в 1924. Два множества $A$ и $B$ в Т. п. наз. ф у кци о на льно о т д е ли м ы ми, если существует определенная на всем пространстве $X$ действительная непрерывная и ограниченная на всем $X$ функция $f$, принимающая во всех точках множества $A$ значение 0 и во всех точках множества $B$ значение 1. В нормальном пространстве всякие два дизъюнктные замкнутые множества функционально отделимы (л е м м а $\mathrm{y}$ р ы с о н а). Обратно, из функциональной отделимости двух любых мно- жеств следует и их отделимость посредством окрестностей. В частности, из функциональной отделимости точки и множества следует их отделимость посредством окрестностей в данном пространстве. Но если пространство регулярно и, значит, всякая точка и всякое не содержащее ее замкнутое множество имеют дизъюнктные окрестности, то отсюда еще не следует, что они функционально отделимы. Таким образом, более сильным, чем свойство регулярности, является свойство по л н о й р е г у л я р н о с т и пространства, заключающееся в том, что всякая точка и всякое не содержащее ее замкнутое множество в этом пространстве функционально отделимы. Среди удовлетворяющих этому условию (или вполне регулярных) пространств наиболее важны вполне регулярные $T_{1}$-пространства, наз. т и х оно в с к и м и п р ос т ра н с т в а и. В частности, пространство всякой топологич. группы является вполне регулярным, но может не быть нормальным.



Наряду с аксиомами отделимости для всей теории Т. п. имеют значение т. н. у с л о в и я т и п а ко мп а к т н о с т и. Они основаны на рассмотрении (открытых) покрытий. Семейство $\Sigma$ (открытых) множеств данного Т. Ш. наз. п о к р ы т и е м әтого пространства $X$, если каждая точка $x \in X$ содержится хотя бы в одном множестве - элементе семейства $\Sigma$. Если каждый әлемент покрытия $\alpha^{\prime}$ является подмножеством хотя бы одного элемента покрытия $\alpha$, то говорят, что $\alpha^{\prime}$ вписано в $\alpha$, или что $\alpha^{\prime}$ мельче $\alpha$, или, наконец, что $\alpha^{\prime}$ следует за $\alpha$ в частично упорядоченном множестве покрытий данного пространства. Частным случаем следования покрытия $\alpha^{\prime}$ за $\alpha$ является случай, когда $\alpha^{\prime}$ содержится в $\alpha$ (т. е. каждый элемент покрытия $\alpha^{\prime}$ есть и элемент покрытия $\alpha$ ). То или иное условие типа компактности предполагает, что даны два класса (открытых) покрытий пространства $X:$ класс $\mathfrak{Y}$ и класс $\mathfrak{B}$ так, что $\mathfrak{B} \subset \mathfrak{A}$, т. е. каждое покрытие класса $\mathfrak{B}$



\includegraphics[max width=\textwidth, center]{2023_07_09_90f4588f098ca4e93dfeg-1(1)}

п а к т н о с т и состоит в том, что за каждым покрытием класса $\mathfrak{U}$ следует покрытие класса $\mathfrak{B}$. Важнейшее среди всех условий этого типа получается, если $\mathfrak{Y}$ есть класс всех открытых покрытий пространства, а $\mathfrak{B}$ - подкласс конечных покрытий, т. е. покрытий, состоящих из конечного числа элементов. Это условие наз. у с ловие м биком па т т о с т и; оно эквивалентно т. н. У с л о в и г Б о р е ля- ЈI е б ег а: для каждого открытого покрытия $\alpha$ пространства $X$ существует конечное покрытие $\alpha^{\prime}$ пространства $X$, содержащееся в $\alpha$. Хаусдорфовы пространства, удовлетворяющие условию бикомпактности, наз. б и к о мп а к т а м и. Все они нормальны. Метризуемые бикомпакты (бикомпакты со счетной базой) наз. к о мп а к т а м и. В настоящее время получила распространение другая терминология, согласно к-рой бикомпакты наз. компактными Т. П., причем метризуемый случай терминологически не выделяется.



Если за $\mathfrak{A}$ принять класс счетных открытых покрытий, сохраняя в качестве $\mathfrak{B}$ подкласс конечных покры-



\includegraphics[max width=\textwidth, center]{2023_07_09_90f4588f098ca4e93dfeg-1}

с ч е т н й ко м п а к т н с т ю ю. Для метризуемых пространств, а также для хаусдорфовых пространств со счетной базой условия бикомпактности и счетной компактности әквивалентны между собой. Если за $\mathfrak{H}$ взять класс всех открытых покрытий, а за $\mathfrak{B}$ - класс счетных покрытий, то получается ус ло в и ф фина льно й ком пактности. При формулировке этого условия можно потребовать (как и в случае бикомпактности), ттобы покрытие $\alpha^{\prime} \in \mathfrak{B}$ не только следовало за покрытием $\alpha$, но и содержалось в нем.



Важный класс пространств, наз. л о к а л ь н о б ик о м ш а к т н м и, определяется требованием, чтобы





\end{document}